% ============================================================================
% PHẦN 3 (2023 - 2024): TỐI ƯU HÓA CHÍNH SÁCH TRỰC TIẾP (DPO)
% ============================================================================
\sectionintro{3}{Direct Preference Optimization (DPO)}{2023--2024}{Closed-form solution, gradient và kiến trúc hai mô hình.}

\subsection{3.1. Từ RLHF Objective đến DPO Loss (Closed-form Solution)}

% ----------------------------------------------------------------------------
% SLIDE: ĐỘNG CƠ CỦA DPO (05/2023)
% ----------------------------------------------------------------------------
\begin{frame}{DPO: tối ưu trực tiếp từ sở thích}
  \lead{Thay thế vòng lặp RL bốn mô hình bằng closed-form solution giải tích (Rafailov et al., 2023).}
  \subhead{Quy trình tối ưu hóa trực tiếp}
  \small
  \begin{itemize}\setlength{\itemsep}{4pt}
    \item \textbf{RLHF + PPO (Phần 2):} Cặp $(y_w, y_l) \xrightarrow{\text{Bước 1}} r_\phi(x, y) \xrightarrow{\text{Bước 2}} \pi_\theta(y \mid x)$ (qua Critic $V_\psi$).
    \item \textbf{DPO (Phần 3):} Biểu diễn giải tích $r(x, y)$ theo Policy $\pi_\theta(y \mid x)$, tối ưu trực tiếp trên cặp $(y_w, y_l)$.
  \end{itemize}

  \vspace{0.10cm}
  \small
  \renewcommand{\arraystretch}{1.15}
  \begin{tabularx}{\textwidth}{@{}l X X@{}}
    \toprule
    \textbf{Tiêu chí} & \textbf{RLHF + PPO (2022)} & \textbf{DPO (2023)} \\
    \midrule
    \textbf{Bộ nhớ VRAM} & \textbf{4 mạng:} $\pi_\theta, \pi_{\text{ref}}, r_\phi, V_\psi$ & \textbf{2 mạng:} Chỉ $\pi_\theta$ và $\pi_{\text{ref}}$ \\
    \textbf{Rollout lúc học} & Bắt buộc ở mọi bước cập nhật & \textbf{Không cần} (Forward song song) \\
    \bottomrule
  \end{tabularx}
  \par\vspace{0pt}
\end{frame}

% ----------------------------------------------------------------------------
% SLIDE: TỪ MỤC TIÊU RLHF ĐẾN NGHIỆM ĐÓNG
% ----------------------------------------------------------------------------
\begin{frame}{DPO: closed-form solution cho RLHF}
  \lead{Bài toán có điều chuẩn KL và nghiệm giải tích Gibbs--Boltzmann.}
  \subhead{Nghiệm cực trị Gibbs--Boltzmann}
  \small
  Xét bài toán cực trị có ràng buộc phân kỳ KL:
  \[
    \fitmath{%
    \max_{\pi} \E_{x,\; y \sim \pi} \left[ r(x, y) - \beta \log \frac{\pi(y \mid x)}{\pi_{\text{ref}}(y \mid x)} \right] \implies \pi^*(y \mid x) = \frac{1}{Z(x)} \pi_{\text{ref}}(y \mid x) \exp\left(\frac{1}{\beta} r(x, y)\right)
    }
  \]
  \vspace{-0.15cm}
  \subhead{Hàm thưởng ngầm định (Implicit Reward)}
  \small
  Nghịch đảo phương trình Gibbs để biểu diễn hàm thưởng theo Policy:
  \[
    \fitmath{%
    \boxed{r(x, y) = \beta \log \frac{\pi^*(y \mid x)}{\pi_{\text{ref}}(y \mid x)} + \beta \log Z(x)}
    }
  \]
  với $Z(x) = \sum_y \pi_{\text{ref}}(y \mid x) \exp\left(\frac{1}{\beta} r(x, y)\right)$ là hàm phân định.
  \par\vspace{0pt}
\end{frame}

% ----------------------------------------------------------------------------
% SLIDE: DPO LOSS
% ----------------------------------------------------------------------------
\begin{frame}{DPO: khử Partition Function và Loss}
  \lead{Triệt tiêu đại lượng $Z(x)$ qua hiệu số cặp Bradley--Terry.}
  \subhead{Khử $Z(x)$ qua hiệu số phần thưởng}
  \small
  Khi lấy hiệu số giữa mẫu dương $y_w$ và mẫu âm $y_l$, $\beta \log Z(x)$ \textbf{tự triệt tiêu}:
  \[
    \fitmath{%
    r(x, y_w) - r(x, y_l) = \beta \log \frac{\pi^*(y_w \mid x)}{\pi_{\text{ref}}(y_w \mid x)} - \beta \log \frac{\pi^*(y_l \mid x)}{\pi_{\text{ref}}(y_l \mid x)}
    }
  \]
  \vspace{-0.15cm}
  \begin{block}{DPO Loss}
    \footnotesize
    Thay $\pi^*$ bằng mạng tham số $\pi_\theta$, tối ưu Maximum Likelihood trên tập cặp:
    \[
      \fitmath{%
      \mathcal{L}_{\text{DPO}}(\theta) = -\E_{(x, y_w, y_l)} \left[ \log \sigma \left( \beta \log \frac{\pi_\theta(y_w \mid x)}{\pi_{\text{ref}}(y_w \mid x)} - \beta \log \frac{\pi_\theta(y_l \mid x)}{\pi_{\text{ref}}(y_l \mid x)} \right) \right]
      }
    \]
  \end{block}
  \par\vspace{0pt}
\end{frame}

\subsection{3.2. Cơ chế Gradient \& Kiến trúc Hai Mô hình}

% ----------------------------------------------------------------------------
% SLIDE: CƠ CHẾ GRADIENT VÀ TỰ ĐIỀU TRỌNG SỐ CỦA DPO
% ----------------------------------------------------------------------------
\begin{frame}{Gradient của DPO}
  \lead{Tự điều chỉnh biên độ cập nhật giữa mẫu dương $y_w$ và mẫu âm $y_l$.}
  \begin{block}{Đạo hàm gradient đầy đủ (Rafailov et al., 2023)}
    \footnotesize
    \[
      \fitmath{%
      \nabla_\theta \mathcal{L}_{\text{DPO}}(\theta) = -\beta \E \Big[ \underbrace{\sigma\big(\hat{r}_\theta(x, y_l) - \hat{r}_\theta(x, y_w)\big)}_{\text{Trọng số lỗi } \in (0, 1)} \cdot \Big( \underbrace{\nabla_\theta \log \pi_\theta(y_w \mid x)}_{\text{Kéo Mẫu dương } (\uparrow)} - \underbrace{\nabla_\theta \log \pi_\theta(y_l \mid x)}_{\text{Đẩy Mẫu âm } (\downarrow)} \Big) \Big]
      }
    \]
  \end{block}

  \subhead{Cơ chế tự điều trọng số động}
  \small
  \begin{itemize}\setlength{\itemsep}{4pt}
    \item \textbf{Khi xếp hạng sai ($\hat{r}_l > \hat{r}_w$):} Hệ số lỗi $\to 1.0 \implies$ Tăng mạnh gradient để sửa sai.
    \item \textbf{Khi phân loại đúng ($\hat{r}_w \gg \hat{r}_l$):} Hệ số lỗi $\to 0 \implies$ Triệt tiêu gradient, chống quá khớp.
    \item \textbf{Tác động đối ngẫu:} Đồng thời tăng $\pi_\theta(y_w)$ ($\uparrow$) và hạ thấp $\pi_\theta(y_l)$ ($\downarrow$).
  \end{itemize}
  \par\vspace{0pt}
\end{frame}

% ----------------------------------------------------------------------------
% SLIDE: SƠ ĐỒ KIẾN TRÚC 2 MÔ HÌNH DPO
% ----------------------------------------------------------------------------
\begin{frame}{DPO: kiến trúc hai mô hình}
  \lead{Chỉ giữ Policy $\pi_\theta$ và Reference Model $\pi_{\mathrm{ref}}$ trên GPU.}
  \vspace{0.15cm}
  \begin{center}
    \resizebox{0.96\textwidth}{!}{%
      \begin{tikzpicture}
        \node[data, text width=3.3cm, minimum height=1.3cm] (data) at (0, 0) {
          \textbf{Cặp đối chiếu tĩnh}\\
          $(x, y_w, y_l) \sim \mathcal{D}_{\text{pair}}$\\
          \normalfont\tiny ($y_w$: Mẫu dương, $y_l$: Mẫu âm)
        };

        \node[trainmodel, text width=3.8cm, minimum height=1.35cm] (actor) at (5.0, 1.2) {
          \textbf{1. Policy Model $\pi_\theta$}\\
          (Đang tối ưu hóa)\\
          \normalfont\scriptsize Tính $\log \pi_\theta(y_w \mid x), \;\log \pi_\theta(y_l \mid x)$
        };

        \node[frozenmodel, text width=3.8cm, minimum height=1.35cm] (ref) at (5.0, -1.2) {
          \textbf{2. Reference Model $\pi_{\text{ref}}$}\\
          (Cố định tham số)\\
          \normalfont\scriptsize Tính $\log \pi_{\text{ref}}(y_w \mid x), \;\log \pi_{\text{ref}}(y_l \mid x)$
        };

        \node[rewardbox, text width=3.7cm, minimum height=2.1cm] (implicit) at (10.1, 0) {
          \textbf{Thưởng Ngầm định}\\[4pt]
          $\hat{r}_w = \beta \log \dfrac{\pi_\theta(y_w \mid x)}{\pi_{\text{ref}}(y_w \mid x)}$\\[4pt]
          $\hat{r}_l = \beta \log \dfrac{\pi_\theta(y_l \mid x)}{\pi_{\text{ref}}(y_l \mid x)}$
        };

        \node[lossbox, text width=3.5cm, minimum height=2.1cm] (dpoloss) at (14.8, 0) {
          \textbf{DPO Loss}\\[4pt]
          $\mathcal{L}_{\text{DPO}} = -\log \sigma(\hat{r}_w - \hat{r}_l)$\\[4pt]
          \normalfont\scriptsize Khử hoàn toàn $r_\phi$\\
          và Value Network $V_\psi$
        };

        \coordinate (bus_in) at (2.25, 0);
        \draw[line width=1.05pt, draw=navyblue!85!black] (data.east) -- (bus_in);
        \draw[flow] (bus_in) |- (actor.west);
        \draw[flow] (bus_in) |- (ref.west);

        \draw[flow] (actor.east) -- (7.55, 1.2) |- ([yshift=0.45cm]implicit.west);
        \draw[flow] (ref.east) -- (7.55, -1.2) |- ([yshift=-0.45cm]implicit.west);
        \draw[flow] (implicit.east) -- (dpoloss.west);

        \draw[gradflow] (dpoloss.north) -- (14.8, 2.45) -- node[above, font=\scriptsize\rmfamily\bfseries, text=stanfordred] {Backpropagation trực tiếp cập nhật $\pi_\theta$ (không cần Rollout)} (5.0, 2.45) -- (actor.north);
      \end{tikzpicture}%
    }
  \end{center}
  \vspace{0.05cm}
  \small
  \textbf{Ưu thế hệ thống:} Tính toán log-xác suất song song trên toàn chuỗi $(y_w, y_l)$ trong một forward pass duy nhất.
  \par\vspace{0pt}
\end{frame}

\subsection{3.3. Tổng kết Phần 3: Hạn chế của DPO trong Suy luận}

% ----------------------------------------------------------------------------
% SLIDE: HẠN CHẾ CỦA DPO
% ----------------------------------------------------------------------------
\begin{frame}{Hạn chế của DPO trong bài toán suy luận}
  \begin{enumerate}\setlength{\itemsep}{4pt}
    \item \textbf{Off-policy Exploration Bottleneck:}\\
      \small Chỉ học thụ động trên tập cặp tĩnh $\mathcal{D}_{\text{pair}}$, không tự khám phá quỹ đạo suy luận mới.
    \item \textbf{Reward Hacking trên bài toán hình thức:}\\
      \small Trong toán và code, mô hình dễ sa vào các lời giải ngụy biện nhưng trình bày đẹp mắt.
    \item \textbf{Không kích hoạt Test-time Scaling:}\\
      \small Không thúc đẩy mô hình mở rộng chuỗi suy nghĩ CoT dài để tự giải quyết bài toán phức tạp.
  \end{enumerate}
  \par\vspace{0pt}
\end{frame}
